\subsection{叉角}

我们也在代数中定义过极大有序域上的二维向量空间中一对直线的叉角（*）。这个定义特别适用于实平面 \(\mathbf{R}^2\) 。全体叉角的集合 \(\mathfrak{A}_0\) 具有一个交换群结构（写成加法），它由下式定义

\[
(\widehat {\mathbf {D _ {1} , D _ {3}}}) = (\widehat {\mathbf {D _ {1} , D _ {2}}}) + (\widehat {\mathbf {D _ {2} , D _ {3}}})
\]

因而特别有 \((\widehat{\mathbf{D}_1,\mathbf{D}_1}) = 0\) 与 \((\widehat{\mathbf{D}_2,\mathbf{D}_1}) = -(\widehat{\mathbf{D}_1,\mathbf{D}_2}).\)

直叉角 \(\delta_0\) 是 \(\neq 0\) 的、方程 \(2\theta = 0\) 在 \(\mathfrak{A}_0\) 中的解；它是虚轴与实轴所成的叉角。

我们定义一个典范同态 \(\varphi\) ：由角群 \(\mathfrak{A}\) 到叉角群 \(\mathfrak{A}_0\) 上，它把射线 \(\Delta\) 与实正半轴所成的角，对应为包含 \(\Delta\) 的直线 D 与实轴所成的叉角。一个叉角 \(\theta_0\) 是 \(\varphi\) 下两个角 \(\theta\) 与 \(\theta + \varpi\) 的像；于是 \(\mathfrak{A}_0\) 同构于 \(\mathfrak{A}\) 对子群 \(\{o, \varpi\}\) 的商群。若在 \(\mathfrak{A}_0\) 上转移商群 \(\mathfrak{A}/\{o, \varpi\}\) 的拓扑（借助与 \(\varphi\) 相伴的双射同态来转移），则 \(\mathfrak{A}_0\) 成为一个紧拓扑群，而 \(\varphi\) 成为 \(\mathfrak{A}\) 到 \(\mathfrak{A}_0\) 上的一个严格态射。

如果把同态 \(\varphi\) （ \(\mathfrak{A}\) 到 \(\mathfrak{A}_0\) 上的）与同态 \(x \to \Im(x/a)\) （ \(\mathbf{R}\) 到 \(\mathfrak{A}\) 上的）复合，就得到一个同态 \(x \to \Im_0(x/a)\) （ \(\mathbf{R}\) 到 \(\mathfrak{A}_0\) 上的）；每个叉角 \(\theta_0 \in \mathfrak{A}_0\) 在这个同态之下对应于实数的一个模 \(\frac{1}{2}a\) 的类，它称为叉角 \(\theta_0\) 的测度（相对于基 \(a\) ）；照惯用说法，这个类中的每个数都称为 \(\theta_0\) 的一个测度，而其中属于区间 \([0, \frac{1}{2}a]\) 的那个称为 \(\theta_0\) 的主测度；叉角 \(\Im_0(x/a)\) 是测度为 \(x\) 的叉角。\(\theta_0\) 的每个测度也是角 \(\theta, \theta + \varpi\) 二者之一的一个测度，这两个角在同态 \(\varphi\) 下的像是 \(\theta_0\) 。

这里同样地，一旦选定了基 a ，当我们谈到一个叉角时，照惯用说法，一般指相对于基 a 的这个叉角的一个测度。

注。我们可以定义 \(\mathbf{C}^*\) 到 \(\mathfrak{A}_0\) 上的一个同态，它把每个复数 \(z \neq 0\) 映为过 \(0\) 与 \(z\) 的直线与实轴所成的叉角。显然这个同态是 \(\varphi\) 与同态 \(z \to \operatorname{Am}(z)\) （ \(\mathbf{C}^*\) 到 \(\mathfrak{A}\) 上的）的复合；因此它是拓扑群 \(\mathbf{C}^*\) 到拓扑群 \(\mathfrak{A}_0\) 上的一个严格态射，而与之相伴的双射同态是商群 \(\mathbf{C}^*/\mathbf{R}^*\) 到 \(\mathfrak{A}_0\) 上的一个同构。

我们知道，若 D 表示一条与实轴成叉角 \(\theta_0\) 的直线，且 \((a, b)\) 是 D 的一对方向比，则叉角 \(\theta_0\) 的正切（记作 tan \(\theta_0\) ）是元素 \(b/a\) （ \(\tilde{\mathbf{R}} (= \infty\) 若 \(a = 0\) ），它也称为直线 D 的斜率。若 \(\theta\) 与 \(\theta + \varpi\) 是在 \(\varphi\) 下以 \(\theta_0\) 为像的两个角，则有 \(\tan \theta_0 = \tan \theta = \tan (\theta + \varpi)\) 。映射 \(\theta_0 \to \tan \theta_0\) 是 \(\mathfrak{A}_0\) 到 \(\tilde{\mathbf{R}}\) 上的一个同胚，因为拓扑空间 \(\mathbf{C}^*/\mathbf{R}^*\) 正是实射影直线 \(\mathbf{P}_1\) ，而且由第六章 § 3 第 3 号，把直线（视为 \(\mathbf{P}_1\) 的点）对应到其斜率的映射是 \(\mathbf{P}_1\) 到 \(\tilde{\mathbf{R}}\) 上的一个同胚。现在，若在 \(\tilde{\mathbf{R}}\) 上转移 \(\mathfrak{A}_0\) 的群结构（借助映射 \(\theta_0 \to \tan \theta_0\) 来转移），我们就在 \(\tilde{\mathbf{R}}\) 上定义了一个交换拓扑群的结构，其中两个元素 \(t_1\) 、\(t_2\) 的积是 \(\frac{t_1 + t_2}{1 - t_1 t_2}\) ，只要 \(t_1\) 、\(t_2\) 属于 R 且 \(t_1 t_2 \neq 1\) ；对不满足这些条件的数对 \((t_1, t_2)\) ，\(t_{1}\) 与 \(t_{2}\) 的积通过把函数 \(\frac{x+y}{1-xy}\) 连续地延拓到 \(\tilde{R}\times\tilde{R}\) 而得到，并且仍记作 \(\frac{t_{1}+t_{2}}{1-t_{1}t_{2}}\) 。

